% =====================================================================
%  03 问题分析
%  每个问题用一段连续文字概述解决思路(不分步骤),
%  只讲"怎么做",不写具体公式、不重述问题内容。
%  分步骤的详细建模见第 4 节,以加粗 Step N 引出。
% =====================================================================
\section{问题分析}

\subsection{问题一的分析}

本文基于传热学与流体力学理论\cite{White2016,Yang2006}建立流固耦合控制
方程组,定义热阻、压降与温度非均匀性三项性能指标;进而提取雷诺数、
普朗特数与结构参数比等主导无量纲量,确认附件数据处于层流工况,为后续
代理模型输入选取提供机理依据;在此基础上,依据达西摩擦、局部阻力叠加
与热阻串联机理构造压降幂律关联式与热阻二次结构关联式,由附件数据最小
二乘标定;最后固定其余参数,定量统计各结构参数对三项指标的影响方向与
幅度,验证机理规律并论证三项指标综合评价的合理性。

\subsection{问题二的分析}

对 84 组样本按 $8:2$ 留出法划分训练/测试集,输入特征做 Z-score 标准化;
考虑到样本量小且映射高度非线性,选用高斯过程回归\cite{Rasmussen2006},
以 Constant$\times$RBF+White 复合核建立多输入--单输出(MISO)映射,
用极大似然估计优化超参数;最后以决定系数、均方根误差与平均绝对百分比
误差评估拟合与泛化能力,确认代理模型可替代 CFD 求解器。

\subsection{问题三的分析}

以 GPR 代理模型为目标函数,在样本支撑域内建立三指标同时最小化的多目标
优化模型(热阻-压降负相关 $r=-0.684$ 表明解集为 Pareto 前沿);采用
NSGA-II\cite{Deb2002} 全局搜索 Pareto 前沿,L-BFGS-B\cite{Byrd1995,Chen2005}
局部精修,并对排数整数约束离散修正;最终以 Min--Max 无量纲化消除量纲
差异,用均等权重理想点法\cite{Hwang1981} 将多目标转化为贴近理想解的
单目标决策,唯一确定综合最优设计方案并回代验证。

\subsection{问题四的分析}

在权重单纯形上采样全场景偏好,建立由偏好权重到最优结构参数的敏感性
映射,量化权重波动对方案的驱动作用;以各场景定制最优为基准定义后悔值,
采用最小最大后悔值准则\cite{Savage1951} 求取最坏偏好下综合性能损失最小
的鲁棒方案;最后以权重空间变异系数与最大后悔值对比鲁棒方案与等权综合
方案,论证鲁棒设计的工程价值。

\subsection{问题五的分析}

对加工误差(结构参数)与工况波动(流量、入口温度、发热功率)两类扰动进行
概率分布表征,建立 $\pm5\%$ 综合摄动模型;结构参数层由代理模型计算局部
归一化敏感性系数与 Sobol 全局敏感性指数\cite{Sobol2001},工况参数层基于
无量纲标度关系一阶近似;采用拉丁超立方采样与蒙特卡洛传播\cite{McKay1979}
生成性能响应分布,以变异系数、95\% 置信上限与越界失效概率判定稳定性;
最后将综合最优方案与鲁棒方案一并代入扰动传播,对比评估两类方案的工程
稳定性。
